% pttai: Build-Time Dataflow Validation for LLM Pipeline Graphs
% EMNLP 2026 System Demonstrations submission.
%
% Style: official ACL/EMNLP style files (acl.sty, acl_natbib.bst) fetched from
%   https://github.com/acl-org/acl-style-files (master).
% Venue is SINGLE-BLIND, so the author block is shown (no [review] anonymization).
% To produce an anonymized, line-numbered review copy instead, change the acl
% option below from `final` to `review`.
% Compile:  tectonic paper/main.tex        (or: latexmk -pdf paper/main.tex)

\documentclass[11pt]{article}

\usepackage[final]{acl}

\usepackage{times}
\usepackage{latexsym}
\usepackage[T1]{fontenc}
\usepackage[utf8]{inputenc}
\usepackage{microtype}
\usepackage{inconsolata}
\usepackage{graphicx}
\usepackage{booktabs}
\usepackage{amssymb}
\usepackage{multirow}
\usepackage{enumitem}

\newcommand{\may}{\textsc{may}}
\newcommand{\must}{\textsc{must}}
\newcommand{\code}[1]{\texttt{#1}}

\title{\code{pttai}: Build-Time Dataflow Validation for LLM Pipeline Graphs}

% TODO(author): fill in the real author block (names, affiliation, email).
\author{
  Anonymous \\
  % TODO(author): Affiliation \\
  % TODO(author): \texttt{email@domain} \\
}

\begin{document}
\maketitle

\begin{abstract}
LLM agent pipelines are graphs of model calls, tools, and branches over a shared
state. A miswired one---a node that reads a state key nothing upstream produced,
an unwired branch, a concurrent write to a key with no reducer---does not fail
when you build it. It fails \emph{after} the model has already run, once
execution reaches the broken node, having burned tokens, latency, and (with
tools) real side effects on a run that was doomed at construction. We present
\code{pttai}, a thin declarative layer in which each node declares the state keys
it \code{reads} and \code{writes}, wired with a \code{>} operator; it compiles to
a native LangGraph~\citep{langgraph} \code{StateGraph} and, \emph{before}
compiling, runs a build-time availability (dataflow) analysis over the graph. The
analysis is a \may/\must\ fixpoint over the graph's acyclic condensation, in the
tradition of reaching-definitions and definite-assignment
analysis~\citep{kildall,dragon,jls}: hard errors are raised only from the \may\
(some-path) set, so it does not reject a pipeline that would have run---we observe
\textbf{0} false positives on 19 valid pipelines, which we report as a
measurement, not a guarantee. On a controlled 17-item bug suite \code{pttai}
fails the build on all 17; on the 12-item subset that raw LangGraph does not also
catch at build, LangGraph surfaces every case only at run time. Our primary,
non-circular evaluation applies the same validator to pipelines a frontier model
generates from NLP task specs---a static guardrail for AI-written agent code
(numbers pending a keyed generation run; \S\ref{sec:eval}). The demonstration is
an interactive playground that paints the offending node red at build time,
before any model call.
\end{abstract}

\section{Introduction}

Language-model pipelines have become graphs. A retrieval-augmented QA system
retrieves passages, then answers grounded in them; an extract-then-summarize
pipeline pulls typed fields from a ticket, then writes a one-line triage summary
from those fields; a document-triage pipeline classifies an incoming message and
routes it to a specialized handler. Frameworks such as
LangGraph~\citep{langgraph} express these as explicit state-machine graphs:
nodes read and write a shared state object, edges (some conditional) move control
between them.

These graphs fail differently from ordinary programs. In a shared-state graph,
a node that reads \code{state["summary"]} before any node has written
\code{summary} does not raise when the graph is built---it raises a
\code{KeyError} at run time, once control reaches that node. By then the pipeline
has already made model calls: tokens spent, latency incurred, and, if earlier
nodes called tools, side effects committed, all on a run that was structurally
broken before it started. The same is true of an unwired conditional branch (a
routing key that resolves to nowhere) and of two parallel branches writing the
same non-reduced state key (a runtime \code{InvalidUpdateError}). These are
recurring, publicly reported failures against raw LangGraph
(\S\ref{sec:wild}).

The root cause is that the graph builder checks \emph{structure} but not
\emph{dataflow}: LangGraph's \code{compile()} does not verify that every state
key a node reads is produced upstream. Ordinary compilers have solved the
analogous problem for decades---definite-assignment analysis rejects a program
that reads a local before it is assigned on some
path~\citep{jls,kildall,dragon}. \code{pttai}'s contribution is not that
algorithm, which is textbook; it is \emph{applying} it to an LLM pipeline's
shared state, and the annotation design (per-node \code{reads}/\code{writes})
that makes the shared-state dataflow explicit enough to analyze at build time.

We present \code{pttai}, a declarative layer over
LangGraph.\footnote{Source, examples, benchmark, and demo:
\url{https://github.com/TeeratP/agentic-framework}.}
Our contributions are:
\begin{enumerate}[noitemsep,topsep=2pt]
  \item \textbf{A declarative \code{>} DSL over LangGraph}
  (\S\ref{sec:design}) whose typed per-node \code{reads}/\code{writes}
  annotations make the shared-state dataflow of a pipeline explicit, while
  compiling down to a native LangGraph \code{StateGraph} so that streaming,
  async, checkpointers, and observability are inherited, not reimplemented.
  \item \textbf{A build-time availability (dataflow) analysis for these graphs}
  (\S\ref{sec:validator}): a forward \may/\must\ fixpoint over the graph's
  acyclic condensation---reaching-definitions / definite-assignment adapted to
  agent graphs~\citep{kildall,dragon,jls}---that fails the build, with zero model
  calls, when a node reads a key no upstream node produces, a branch is unwired,
  or parallel branches write an unreduced key. Because hard errors come only from
  the \may\ (some-path) set, we observe no false positives on our valid corpus.
  \item \textbf{An evaluation} (\S\ref{sec:eval}) whose primary, non-circular
  arm applies the validator to \emph{LLM-generated} pipelines, supported by a
  controlled bug suite and a survey of wild bugs.
\end{enumerate}
\code{pttai} borrows the ``ergonomic layer over a powerful runtime'' idea from
Keras~\citep{keras}, but the ergonomics are not the contribution; the validator
is.

\begin{figure}[t]
  \centering
  \includegraphics[width=\columnwidth]{figures/architecture.png}
  \caption{The \code{>} DSL builds a linked node structure in memory;
  \code{AgenticGraph} walks it into \code{add\_node}/\code{add\_edge}/\code{Send}
  calls; the build-time dataflow validator gates the build
  (\code{GraphValidationError} on failure, \emph{before any model call}); on
  success it compiles to a native LangGraph \code{StateGraph}.}
  \label{fig:arch}
\end{figure}

\section{System Design}
\label{sec:design}

\subsection{The \code{>} DSL and compilation}

A \code{pttai} pipeline is built by instantiating \code{Node} objects and wiring
them with \code{>} (Figure~\ref{fig:arch}). Wiring is deferred: \code{a > b} sets
\code{a.child = b} and returns \code{b}, so \code{a > b > c} builds a linked
structure in memory---no graph edges yet. \code{AgenticGraph(start, end\_nodes)}
then walks that structure and emits the actual LangGraph
\code{add\_node}/\code{add\_edge}/\code{add\_conditional\_edges} calls before
compiling. There are three node types:
\code{AgentNode} (a system-prompted model call with an internal tool-call loop),
\code{DecisionNode} (constrained structured-output routing: the model must
return one of the declared \code{choices}, wired \code{decision["c"] > handler}),
and \code{InputNode} (human-in-the-loop via LangGraph's \code{interrupt}). Fan-out
is \code{fanout(a, b)} and map-reduce is \code{.map(field)}.

Nodes read and write a shared, reducer-based state. Each node declares its
\code{reads} and \code{writes}; a node returns only the keys it updates and the
channel reducers merge them, which is what makes parallel branches and
checkpointing correct. Those declarations are also what the validator analyzes.
Crucially, an \code{AgenticGraph} \emph{is} a LangGraph \code{StateGraph} and its
compiled artifact is LangGraph's \code{CompiledStateGraph}: \code{pttai} adds a
build pass on top of LangGraph and then hands the compiled graph to it. Nothing
in the runtime is reimplemented.

\subsection{The build-time dataflow validator}
\label{sec:validator}

When \code{AgenticGraph(\dots)} is constructed, it runs a static analysis over
the wired nodes and \textbf{fails the build} on a hard error, raising
\code{GraphValidationError} before \code{compile()}. This is the contribution.

\paragraph{The analysis.} During the build we record a tagged edge list:
each edge is \code{(src, dst, kind)} where \code{kind} is \code{seq} (an
unconditional, \textsc{and} edge) or \code{cond} (an exclusive, \textsc{or} edge:
a \code{DecisionNode} choice or a \code{Send} fan-out target). We compute the
graph's acyclic condensation by excluding loop-back edges (found by a DFS), then
compute, per node, two sets of available state keys by forward fixpoint:
\begin{itemize}[noitemsep,topsep=2pt]
  \item \may---keys available on \emph{some} path into the node. A
  monotone-\emph{increasing} union fixpoint.
  \item \must---keys guaranteed on \emph{every} path into the node. A
  monotone-\emph{decreasing} fixpoint.
\end{itemize}
The transfer function is the standard forward one---a node's out-set is its
in-set plus the keys it writes---with two combining rules at joins that
distinguish this from a scalar CFG analysis:
\emph{union across an \textsc{and}-parallel join} (a node after \code{fanout(a,b)}
sees both branches' writes), and \emph{intersect across an exclusive
\textsc{or}-group} (a node after a \code{DecisionNode} is guaranteed only the keys
every choice branch produces). The seed is the \emph{input} keys---reduced
channels, schema keys no node writes, and user-declared inputs---not all declared
keys; seeding with inputs rather than the whole schema is what lets the analysis
catch a read of a \emph{computed} key before its producer has run. Excluding
loop-back edges makes the analysis reflect first-iteration reality: a key
produced only by a downstream node that cycles back is \emph{not} credited as
available on entry, so loop-carried read-before-write is flagged rather than
silently accepted. The fixpoint iterates until a full sweep changes nothing,
bounded by the node count.

\paragraph{Errors from \may, warnings from \must.} Hard \emph{errors} are raised
only from the \may\ set: a read is an error only when the key is available on
\emph{no} path into the node. If a key is available on some but not all paths
(\may\ but not \must), that is a \emph{warning}, never a build failure. This is
why we observe no false positives: the analysis reports ``used before produced''
only when \emph{no} path produces the key, exactly as definite-assignment
analysis does~\citep{jls,nielson}. We are careful \emph{not} to claim this as a
soundness proof (see Limitations): it is sound for the errors it raises given the
declared annotations, not complete, and the 0-false-positive figure is a
measurement on our corpus.

\paragraph{What it checks.} On top of the dataflow sets, \code{pttai} raises
seven error classes and four warning classes (Table~\ref{tab:checks}). Errors
fail the build; warnings surface in a \code{model.summary()}-style table and can
be promoted to errors with \code{strict=True}.

\begin{table}[t]
  \centering
  \small
  \begin{tabular}{@{}p{0.42\columnwidth}p{0.5\columnwidth}@{}}
    \toprule
    \textbf{Error (fails build)} & \textbf{Condition} \\
    \midrule
    read-before-write & reads a computed key produced only downstream / on a
    sibling branch (\may) \\
    undeclared key & reads/writes a key absent from the state schema \\
    concurrent write, no reducer & two parallel branches write the same
    non-reduced key \\
    dangling choice & a \code{DecisionNode} choice with no wired handler \\
    dead-end node & a non-terminal node with no child$^\ddagger$ \\
    duplicate node name & two distinct nodes share a name$^\dagger$ \\
    prompt/read mismatch & a \code{node\_prompt} \code{\{placeholder\}} that is
    not a declared scalar read \\
    \midrule
    \textbf{Warning (build succeeds)} & \textbf{Condition} \\
    \midrule
    some-paths read & key in \may\ but not \must \\
    optional read unmet & \code{reads=["x?"]} never produced \\
    unreachable node & not reachable from the start node \\
    end node w/ children & a declared end node still has a child \\
    \bottomrule
  \end{tabular}
  \caption{The validator's error and warning classes. $^\dagger$Duplicate names
  and dead ends are raised structurally during the build walk (as
  \code{ValueError}); the rest come from the dataflow pass. $^\ddagger$The
  dead-end rejection is a \code{pttai} DSL-strictness choice, \emph{not} a
  LangGraph bug: LangGraph treats a childless node as a legal implicit terminal
  (\S\ref{sec:bugbench}). A dead read (a scalar read never interpolated) is an
  additional warning.}
  \label{tab:checks}
\end{table}

\section{Demonstration}
\label{sec:demo}

The demonstration is a browser playground (Figure~\ref{fig:demo}). A reviewer
edits a small \code{pttai} snippet---one of the shipped NLP pipelines (RAG QA,
extract$\rightarrow$summarize, document triage) or their own---and submits it. On
a clean build the playground renders the \emph{compiled LangGraph} as a
node/edge diagram plus the validator's \code{summary()} table (every node's
reads, writes, and available keys). The intended live moment: take a working
RAG-QA pipeline and reorder two nodes so the answer step reads \code{passages}
before the retriever runs. Raw LangGraph would compile this and only fail after
the retrieval model call. \code{pttai} instead rejects it at build: the diagram
is redrawn from the pre-compile wiring with the \textbf{offending node painted
red} and the exact error attached (``\code{reads computed key 'passages' but no
upstream node produces it}''), with \emph{zero model calls made}. The reviewer
sees the picture of the bug before running anything.

The playground itself never invokes a model---it builds, validates, and draws.
To show that \code{pttai} \emph{is} a working agent framework and not only a
linter, the accompanying screencast first runs a live pipeline (the
parallelization example, which really calls the model) end-to-end, and only then
switches to the playground for the build-time catch. Because the playground
\code{exec}s user-supplied code, it is sandboxed for public hosting: the snippet
is size-capped, AST-checked (imports allow-listed to
\code{pttai}/\code{typing}/safe stdlib; dunder access and dangerous builtins such
as \code{eval}/\code{open}/\code{\_\_import\_\_} rejected), run with a restricted
\code{\_\_builtins\_\_}, and executed under a per-request timeout. The model
backend is an offline fake, so the demo needs no API key and makes no network
call. A \textless2.5-minute screencast accompanies the submission.

\begin{figure}[t]
  \centering
  \includegraphics[width=\columnwidth]{figures/demo_screenshot.png}
  \caption{The playground. A broken pipeline is rejected at build time and drawn
  with the offending node in red and the validator error attached---no model
  call is made.}
  \label{fig:demo}
\end{figure}

\section{Evaluation}
\label{sec:eval}

All evaluations run offline (no API key); the harness is in the repository under
\code{eval/}. The controlled suite and wild-bug survey are authored by us and are
therefore vulnerable to the objection that a validator catches only the bugs its
author planted. Our \emph{primary} evaluation is designed to remove that
objection.

\subsection{Primary: LLM-generated pipelines}
\label{sec:llmgen}

The most 2026-relevant question is whether the validator is a useful static
guardrail for \emph{AI-written} agent code. We therefore designed a study
(\code{eval/llmgen/}) in which a frontier model generates \code{pttai} pipelines
from short NLP task specs---given only the docs, and \textbf{never told to write
bugs}---after which every generated pipeline is built through the validator and
every flag is author-adjudicated as a true bug or a false positive per a fixed
protocol (\code{adjudicate.md}). Because the bug distribution comes from a model
solving ordinary tasks, whatever the validator catches, misses, or wrongly flags
is a fact about real AI-written code, not about a hand-picked corpus. The
protocol targets $N\!\ge\!100$ generated pipelines; construction alone scores a
pipeline (the validator runs inside \code{AgenticGraph.\_\_init\_\_}), so scoring
needs no model call, only the generation step does. The scoring path is verified
offline on a committed 9-pipeline sample set (4 clean, 5 flagged, 0 false
positives), but that is a harness self-test, not the study.

\begin{quote}\small
\textbf{TODO(author):} populate from \code{eval/llmgen/run\_study.sh}. Report
$N$, the fraction of \emph{buildable} pipelines flagged, the
\emph{author-adjudicated} false-positive rate among flags, and the split of flags
by the LangGraph phase (runtime / silent / also-at-build). These require an
\code{OPENAI\_API\_KEY} and a keyed generation run that the build environment
cannot perform; \textbf{we do not report or claim them here}.
\end{quote}

\subsection{Controlled bug suite}
\label{sec:bugbench}

\code{bugbench} is a labeled unit-test suite of 17 buggy and 19 valid pipelines
that pins the validator's behavior per bug class; it is not the headline result.
For each item we record whether \code{pttai} fails the build and---using an
offline fake model---in which phase raw LangGraph would surface the same bug.
\code{pttai} fails the build on \textbf{all 17} buggy pipelines with the correct
diagnostic (median build time under 1\,ms) and \textbf{0} false positives on the
19 valid pipelines. We partition the buggy set into three honest categories
(Table~\ref{tab:bugbench}):

\begin{itemize}[noitemsep,topsep=2pt]
  \item \textbf{Differentiator (12 items, 5 classes)}: read-before-write,
  cyclic-read-before-write, dangling-choice, concurrent-write-no-reducer,
  prompt-placeholder-mismatch. \code{pttai} catches all 12 at build; raw
  LangGraph catches \textbf{0} at build and surfaces all \textbf{12} at run
  time.
  \item \textbf{Caught by both (2 items)}: duplicate node names. Both frameworks
  reject these at build (\code{ValueError}), so they are \emph{not} a
  differentiator---we say so explicitly.
  \item \textbf{DSL-strictness (3 items, excluded)}: dead-end nodes. LangGraph
  1.2 treats a childless node as a legal implicit terminal that compiles and runs
  (we verified this empirically); \code{pttai} rejecting it is a strictness
  choice of our DSL, \emph{not} a LangGraph bug, so we exclude these from the
  differentiator claim.
\end{itemize}

On the differentiator subset, the offline harness attributes \textbf{8}
LangGraph model calls to reaching the bug against \code{pttai}'s \textbf{0}. We
flag this number as \emph{simulated}: it is produced by a deterministic offline
fake model under a worst-case node ordering, not a measurement of real token cost.

\begin{table}[t]
  \centering
  \small
  \begin{tabular}{@{}lccc@{}}
    \toprule
    \textbf{Category / class} & \textbf{$n$} & \textbf{pttai} &
    \textbf{LangGraph} \\
    \midrule
    \multicolumn{4}{@{}l}{\emph{Differentiator (build-catch only in pttai)}} \\
    read-before-write        & 3 & build & runtime \\
    cyclic-read-before-write & 2 & build & runtime \\
    dangling-choice          & 3 & build & runtime \\
    concurrent-write/no-red. & 2 & build & runtime \\
    prompt/read mismatch     & 2 & build & runtime \\
    \emph{subtotal}          & \textbf{12} & \textbf{12/12} & \textbf{0/12} \\
    \midrule
    \multicolumn{4}{@{}l}{\emph{Caught by both (not a differentiator)}} \\
    duplicate node names     & 2 & build & build \\
    \midrule
    \multicolumn{4}{@{}l}{\emph{DSL-strictness (excluded)}} \\
    dead-end node            & 3 & build & runs \\
    \bottomrule
  \end{tabular}
  \caption{\code{bugbench} outcomes by category (\code{eval/bugbench/results.json}).
  ``build'' = rejected at construction; ``runtime'' = raises only at
  \code{invoke()}; ``runs'' = LangGraph executes it as a legal terminal.
  \code{pttai} additionally raises 0 false positives on 19 valid pipelines.}
  \label{tab:bugbench}
\end{table}

\subsection{Wild-bug survey}
\label{sec:wild}
To show the taxonomy is not author-imagined, we mined public issue trackers
(Table~\ref{tab:wild}): six independent, verified reports against raw LangGraph,
two per class, each a runtime failure that \code{pttai}'s build-time check would
have caught first. This is an \emph{illustrative sample}, not a frequency claim:
these are six issues drawn from thousands, with no denominator, so they establish
that the classes occur in the wild, not how often. (Duplicate names and the
prompt/read check are excluded: the former both frameworks already catch at
build, and the latter is a \code{pttai}-internal reformulation of a generic
\code{str.format} error with no LangGraph-specific report.)

\begin{table}[t]
  \centering
  \small
  \begin{tabular}{@{}p{0.6\columnwidth}c@{}}
    \toprule
    \textbf{Bug class (runtime symptom)} & \textbf{Reports} \\
    \midrule
    read-before-write (\code{KeyError} on a state key) & 2 \\
    concurrent write, no reducer (\code{InvalidUpdateError}) & 2 \\
    dangling conditional edge & 2 \\
    \midrule
    \textbf{Total} & \textbf{6} \\
    \bottomrule
  \end{tabular}
  \caption{Independent, publicly reported wild bugs against raw LangGraph, by
  targeted class---an illustrative sample, not a census (details and issue links
  in \code{docs/EVIDENCE.md}).}
  \label{tab:wild}
\end{table}

\subsection{Conciseness and overhead}
As a side effect of folding \code{add\_node}/\code{add\_edge}/structured-output
routing/the tool-call loop into the node abstractions, the \code{pttai} versions
of 12 canonical pipelines total 113 lines against 281 for hand-written LangGraph
(\code{eval/loc\_results.csv}); this is a consequence of the DSL, not the point of
the system. The added build pass has a sub-millisecond to low-single-digit-ms
build time on every corpus pipeline (\code{eval/bugbench/results.json}); because
\code{pttai} compiles to a native \code{CompiledStateGraph}, run time, model-call
count, and output are LangGraph's.

\section{Related Work}
\label{sec:related}

\code{pttai} is complementary to, not a replacement for, existing frameworks; it
adds a build-time state-key dataflow check and compiles down to LangGraph.
\textbf{LangGraph}~\citep{langgraph} is the runtime \code{pttai} targets; its
\code{compile()} runs basic structural checks but does not analyze which state
keys are produced before they are read---the gap this work fills. Several
frameworks already perform genuine build-time \emph{structural} checks, so we do
\emph{not} claim to be the first build-time check of any kind:
\textbf{pydantic-graph} (PydanticAI)~\citep{pydanticgraph} encodes a node's
out-edges in the return-type annotation of its \code{run} method, so a standard
static type checker (mypy/pyright) flags a mis-wired edge before the graph runs---
a build-time graph-\emph{structure} check---but its state is a single shared
object passed uniformly to every node, with no per-node read/write availability
analysis; \textbf{AutoGen}'s GraphFlow~\citep{autogen} and
\textbf{LlamaIndex} Workflows~\citep{llamaindex} validate graph structure (entry
points, reachability, event-type wiring), and \textbf{Haystack}~\citep{haystack}
type-checks explicit component connections. What none of these does is a
shared-state \emph{dataflow} availability analysis (read-before-write /
concurrent-unreduced-write). \textbf{DSPy}~\citep{dspy} optimizes prompts against
a metric; its control flow is arbitrary Python, so there is no declared graph to
check. \textbf{LMQL}~\citep{lmql} and \textbf{Guidance}~\citep{guidance} constrain
\emph{within} a single generation rather than across a multi-node graph, and
\textbf{CrewAI}~\citep{crewai} orchestrates via roles with no declared node/edge
graph. To our knowledge, among the frameworks we surveyed \code{pttai} is the
only one that runs a build-time state-key \emph{dataflow} availability check, and
the only one that lowers to an established graph runtime rather than shipping its
own.

\section*{Limitations}

We state the analysis's holes as owned limitations. (1) The \may/\must\ analysis
is \emph{sound for the errors it raises but not complete}: it is not a soundness
proof, and it reasons over \emph{availability}, not types or values---a key
produced with the wrong type or an empty value still passes. (2) It analyzes only
the \emph{declared} \code{reads}/\code{writes} and the graph edges; it does
\textbf{not} read the code inside \code{ConditionNode} predicates (arbitrary
lambdas), tool bodies, or other custom callables, so a state key touched only
inside such a callable is invisible unless the node declares it. (3) After the
loop-back-edge fix the analysis catches loop-carried read-before-write on the
first iteration, but it makes no termination or runaway-loop guarantee: a graph
that loops forever still passes the build. (4) The 0-false-positive figure is
\emph{observed on 19 valid pipelines}, not a guarantee; the \must\ set can also
over-approximate exclusivity where a decision handler merges back through an
unconditional edge, which can only \emph{under}-warn, never produce a wrong hard
error. (5) Requiring \code{reads=}/\code{writes=} declarations means an undeclared
but genuinely-produced input can be rejected until it is declared, and the DSL
itself has wiring footguns (e.g.\ \code{[b,c] > [d,e]}). (6) The wild-bug survey
is an illustrative sample, not a census; and the primary LLM-generation study is
designed and its scoring path verified, but its headline numbers await a keyed
generation run (\S\ref{sec:eval})---we do not report or claim them here.

\section*{Ethics and Broader Impact}

\code{pttai}'s goal is to catch broken agent pipelines \emph{before} they run,
which directly reduces wasted model calls---and thus the tokens, cost, latency,
and energy spent on runs that were doomed at build time---and avoids committing
tool side effects on such runs. The public playground \code{exec}s
user-submitted code; we sandbox it (allow-listed imports, restricted builtins,
AST rejection of dangerous calls, size and time caps, offline fake model) so that
hosting it does not expose the host, but we note that \code{exec}-based sandboxes
are best-effort and a hardened deployment should add OS-level isolation. The
validator makes no claim about the \emph{safety, factuality, or bias of model
outputs}; it checks pipeline dataflow, not what the models say. \code{pttai} is
released under the MIT license.

\bibliography{refs}

\appendix

\section{Reproducing the evaluation}
\label{app:repro}

All controlled numbers are produced offline (no API key) from the repository:
\code{eval/bugbench/run.py} regenerates \code{eval/bugbench/results.json} (the 17
buggy / 19 valid outcomes, per-class categories, and simulated wasted-call
counts); \code{eval/loc\_compare.py} regenerates \code{eval/loc\_results.csv} (the
113 vs.\ 281 line counts over 12 pipelines); \code{eval/overhead.py} prints the
build-time and runtime-parity measurements; and
\code{eval/llmgen/score.py --gen-dir samples} runs the offline sample-set
self-test of the primary study. The primary study's headline numbers require an
\code{OPENAI\_API\_KEY} and \code{eval/llmgen/run\_study.sh}. The wild-bug reports
with issue links are in \code{docs/EVIDENCE.md}.

\section{Example: a validated NLP pipeline}
\label{app:example}

The extract-then-summarize pipeline. The \code{extract} node returns typed
structured output; the \code{summarize} node reads those scalar fields and the
validator confirms, at build time, that each is produced upstream before it is
read.

\begin{quote}\small\ttfamily
extract = AgentNode(llm=llm,\\
\ \ node\_prompt="Extract fields: \{ticket\}",\\
\ \ reads=["ticket"],\\
\ \ writes=\{"product": str, "severity": int\})\\
summarize = AgentNode(llm=llm,\\
\ \ node\_prompt="severity-\{severity\} on \{product\}",\\
\ \ reads=["product", "severity"],\\
\ \ writes=["summary"], output\_field="summary")\\
extract > summarize\\
AgenticGraph(start\_node=extract,\\
\ \ end\_nodes=\{summarize\})
\end{quote}

Swapping the two nodes (\code{summarize > extract}) makes \code{summarize} read
\code{product} and \code{severity} before \code{extract} produces them; the build
fails with a read-before-write error and no model call is made.

\end{document}
